% TOP_PROBLEM: {"id": 5300008, "title": "Limit of tunings along growing periods", "category_id": 11, "category": {"id": 11, "name": "dynamical_systems", "display_name": "Dynamical Systems", "description": "Problems about long-term behavior of deterministic systems, Hamiltonian dynamics, and periodic orbits.", "slug": "dynamical-systems", "order_index": 11, "created_at": "2026-07-31T15:26:25.671Z"}, "status": "open"}
% FIRST_POSTED: null
% ENTRY_KIND: statement_only
% Imported from: batch04.tex; UnsolvedMath ID 5300008
\documentclass[11pt]{article}
\usepackage{fontspec}
\setmainfont{Latin Modern Roman}
\setsansfont{Latin Modern Sans}
\usepackage{amsmath,amssymb,mathtools}
\usepackage{unicode-math}
\setmathfont{Latin Modern Math}
\usepackage{xurl}
\usepackage[hidelinks]{hyperref}
\newcommand{\sref}[2]{\hyperlink{source:#1:#2}{[S#2]}}
\newcommand{\cataloguescope}[1]{\par\noindent\textbf{Question.} #1\par}
\begin{document}
% UnsolvedMath ID 5300008; source catalogue code AMR-052-0008
\setcounter{section}{5300007}
\section{Limit of tunings along growing periods}\label{5300008}
{\small\sffamily UnsolvedMath ID 5300008 · Dynamical Systems}
\subsection{Definitions and mathematical statement}

\paragraph{Shared notation.}$\mathbb N=\{1,2,\ldots\}$, and $\mathbb Z,\mathbb Q,\mathbb R,\mathbb C$ denote the integers, rationals, reals and complex numbers. Any other conventions are specified locally.


For a complex polynomial $P$ of degree at least two, its filled Julia set $K(P)$ is the set of points with bounded forward orbit, and $J(P)=\partial K(P)$ is its Julia set. An external ray is a radial curve in the coordinate near infinity that conjugates $P$ to its leading power map, continued into the basin of infinity. Its angle is measured in $\mathbb R/\mathbb Z$. A critical point is a zero of $P'$. Topological conjugacy means a homeomorphism $h$ satisfying $h\circ f=g\circ h$. A quasiconformal homeomorphism has bounded infinitesimal eccentricity; quasiconformal surgery cuts and glues dynamical planes and finds an invariant complex structure that can be straightened to a holomorphic map.
Let $P_1$ have a critical point $\omega$ of exact period $k\geq1$, and suppose that no other critical point lies anywhere in the full attracting basin of that superattracting cycle. Let $d\geq2$ be the local degree of $P_1$ at $\omega$. Let $P_2$ have exactly one finite critical point and degree $d$. Require $J(P_1)$ and $J(P_2)$ to be connected. Let $B$ be the immediate basin of $\omega$, assume $\overline B$ is homeomorphic to a closed disk, and assume $J(P_2)$ is locally connected. The return map $P_1^k$ on $\overline B$ is conjugated to $z\mapsto z^d$ on the closed unit disk. If $d>2$, choose one of the $d-1$ possible conjugating homeomorphisms; this choice fixes the internal-angle marking.
Replace $B$ by the dynamical plane of $P_2$ compactified by its circle at infinity, matching external angles with these internal angles, and collapse the external rays. Make the corresponding replacement at every inverse image of $B$. The dynamics on the inserted copy in $B$ is given by $P_2$; between the other inserted inverse-image copies it is the identity in the corresponding coordinates, and elsewhere it retains the host dynamics. Denote the resulting topological map by $\operatorname{Tune}(P_1,P_2;\omega)$. The expected polynomial realization has degree $\deg P_1$. The source also includes the quadratic tuning construction when $J(P_2)$ is not locally connected; that quadratic extension is included in the questions below.
In questions about variation, polynomial coefficients carry their usual topology; on a fixed-degree space this is equivalent to locally uniform convergence. Rational maps are considered up to conformal conjugacy, with the topology inherited from their fixed-degree coefficient space. The identifications used in the surgery are kept as part of the construction when the inputs vary.
\paragraph{Statement recorded in the supplied source.}
Let $P_{1,j}$ be admissible hosts with marked superstable-cycle periods tending to infinity, and suppose $P_{1,j}\to P_{1,\infty}$. For a fixed admissible $P_2$, do their polynomial tunings also converge to the same limit $P_{1,\infty}$? \sref{5300008}{2}
\subsection{Short English statement}
When host polynomials converge while their superstable periods grow without bound, do their tunings by a fixed polynomial converge to that same host limit?
\subsection{Sources}
\begin{description}
\item[\hypertarget{source:5300008:1}{S1}] ULAM.AI and the UnsolvedMath Contributors, UnsolvedMath: A Curated Collection of Open Mathematics Problems, record 5300008 (AMR-052-0008). CC BY 4.0. \url{https://huggingface.co/datasets/ulamai/UnsolvedMath}.
\item[\hypertarget{source:5300008:2}{S2}] Ben Bielefeld, Questions in Quasiconformal Surgery, in Ben Bielefeld and Mikhail Lyubich (editors), Problems in holomorphic dynamics (1992). Tuning definition and Question 8, pp. 3–4. \url{https://arxiv.org/pdf/math/9205209}.
\end{description}
\label{end:5300008}
\end{document}
